\documentclass[a4paper,11pt]{article}

\usepackage[utf8]{inputenc}
\usepackage[T1]{fontenc}
\usepackage[ngerman]{babel}
\usepackage{amsmath}
\usepackage{amssymb}
\usepackage{xcolor}
\usepackage[most]{tcolorbox}
\usepackage{titlesec}
\usepackage{enumitem}
\usepackage{tikz}
\usetikzlibrary{arrows.meta}
\usepackage{geometry}
\geometry{a4paper,margin=2.0cm}

% ===== Kitty-Theme (serenity) =====
\definecolor{bgdark}{HTML}{11121A}
\definecolor{panel}{HTML}{1B1F2C}
\definecolor{fgtext}{HTML}{D5DDE8}
\definecolor{accent}{HTML}{91A8D0}
\definecolor{accentlt}{HTML}{B3CEE5}
\definecolor{green}{HTML}{B5D4A8}
\definecolor{yellow}{HTML}{E8D49A}
\definecolor{rose}{HTML}{D8909A}
\definecolor{purple}{HTML}{B69ACE}
\definecolor{cyan}{HTML}{8FBED1}
\definecolor{muted}{HTML}{8F9BB3}
\definecolor{rulecol}{HTML}{4A5B78}
\pagecolor{bgdark}
\color{fgtext}
\renewcommand{\normalcolor}{\color{fgtext}}
\usepackage[colorlinks=true,linkcolor=accentlt,urlcolor=accentlt]{hyperref}

\setlength{\parindent}{0pt}
\setlength{\parskip}{4pt}

\titleformat{\section}{\Large\bfseries\color{accent}}{\thesection}{0.5em}{}
\titleformat{\subsection}{\large\bfseries\color{cyan}}{\thesubsection}{0.5em}{}
\titlespacing*{\section}{0pt}{12pt}{3pt}
\titlespacing*{\subsection}{0pt}{7pt}{2pt}

% Status-Chip
\newcommand{\chip}[2]{{\setlength{\fboxsep}{2.5pt}\colorbox{#1}{\color{bgdark}\footnotesize\bfseries #2}}}
\newcommand{\stark}{\chip{green}{stark}}
\newcommand{\teils}{\chip{yellow}{teils}}
\newcommand{\ueben}{\chip{rose}{üben!}}
\newcommand{\offen}{\chip{muted}{ungetestet}}

% genau EIN Kasten-Typ: durchgerechnetes Beispiel
\newtcolorbox{rechnung}[1]{enhanced,breakable,colback=panel,coltext=fgtext,boxrule=0.5pt,
  colframe=rulecol,coltitle=accentlt,colbacktitle=panel,fonttitle=\bfseries\small,
  title={Durchgerechnet -- #1},arc=2pt,left=5pt,right=5pt,top=3pt,bottom=3pt,
  before skip=5pt,after skip=6pt}

% TikZ-Defaults
\tikzset{ax/.style={rulecol,-{Stealth[length=4pt]}},
         crv/.style={accent,very thick},
         hl/.style={rose,very thick}}

\setlist[itemize]{leftmargin=1.4em,itemsep=1pt,topsep=2pt}
\setlist[enumerate]{leftmargin=1.7em,itemsep=1pt,topsep=2pt}

\begin{document}

\begin{center}
{\Huge\bfseries\color{accent} Numerische Methoden}\\[3pt]
{\large\color{fgtext} Komplettes Lernskript -- allgemein, mit Graphen, für jede Klausur}\\[3pt]
{\small\color{muted} Erklärt \emph{was} es ist, \emph{wie} es geht und \emph{wann} man es nutzt -- pro Thema ein durchgerechnetes Beispiel.}
\end{center}
\vspace{2pt}
\noindent\textcolor{rulecol}{\rule{\linewidth}{0.4pt}}

{\small\color{muted}Dein Stand pro Thema (aus der letzten Probeklausur):\quad
\stark\ = im Test gut \quad \teils\ = kleine Lücke \quad \ueben\ = war Lücke \quad
\offen\ = noch nicht geprüft}

{\color{muted}\tableofcontents}
\noindent\textcolor{rulecol}{\rule{\linewidth}{0.4pt}}

% =====================================================================
\section{Zahlendarstellung \& Fehler}
\stark

\subsection{Warum überhaupt Fehler?}
Ein Rechner hat \emph{endlich} viele Bits. Reelle Zahlen müssen also \emph{gerundet} werden, und
viele Verfahren brechen einen Grenzprozess ab. Jedes numerische Ergebnis trägt deshalb einen Fehler.
Man unterscheidet drei Quellen:
\begin{itemize}
\item \textbf{Modellfehler:} die Realität wird vereinfacht (z.\,B.\ Reibung ignoriert).
\item \textbf{Abschneide-/Diskretisierungsfehler:} ein Grenzwert wird abgebrochen (Reihe, Ableitung
  $\to$ Differenzenquotient). Skaliert mit der Schrittweite $h$.
\item \textbf{Rundungsfehler:} endliche Stellenzahl der Maschine.
\end{itemize}
\textbf{Absoluter} Fehler $|x-\tilde x|$, \textbf{relativer} Fehler $\frac{|x-\tilde x|}{|x|}$
(skaleninvariant -- „wie viele gültige Stellen``; bei $x\approx 0$ nur absolut sinnvoll).

\subsection{Gleitkomma, Maschinenepsilon, Auslöschung}
Eine Gleitkommazahl ist $x=(-1)^s\,m\cdot 2^{e}$: Vorzeichen $s$, \emph{Mantisse} $m$ (signifikante
Stellen $\to$ \textbf{Auflösung}), \emph{Exponent} $e$ (Skalierung $\to$ \textbf{Wertebereich}).
Normalisiert heißt $m=1{.}b_1b_2\dots$ mit implizitem führenden Bit. Die zentrale Kennzahl ist das
\textbf{Maschinenepsilon} $\varepsilon=2^{-t}$ ($t$ = Mantissen-Nachkommabits): der Abstand von $1$
zur nächsten darstellbaren Zahl.

Wichtig ist das Bild der \emph{Abstände}: darstellbare Zahlen liegen nicht gleichmäßig, sondern
werden mit jeder Zweierpotenz \emph{doppelt so weit} auseinander -- konstante \emph{relative},
wachsende \emph{absolute} Auflösung:

\begin{center}
\begin{tikzpicture}[scale=1.15]
  \draw[ax] (0.6,0) -- (8.7,0) node[right]{$x$};
  \foreach \v in {1,2,4,8} \draw[rulecol](\v,-0.07)--(\v,0.07) node[below=2pt]{\scriptsize \v};
  % [1,2): Abstand 0,25
  \foreach \v in {1,1.25,1.5,1.75} \fill[green](\v,0) circle(1.3pt);
  % [2,4): Abstand 0,5
  \foreach \v in {2,2.5,3,3.5} \fill[yellow](\v,0) circle(1.3pt);
  % [4,8): Abstand 1
  \foreach \v in {4,5,6,7} \fill[rose](\v,0) circle(1.3pt);
  \node[green]  at (1.375,0.42){\scriptsize Abst.\ $0{,}25$};
  \node[yellow] at (2.75,0.42){\scriptsize Abst.\ $0{,}5$};
  \node[rose]   at (5.5,0.42){\scriptsize Abst.\ $1$};
\end{tikzpicture}\\[-2pt]
{\footnotesize\color{muted}Gleich viele Zahlen pro „Oktave`` $[1,2),[2,4),[4,8)$ -- der absolute Abstand verdoppelt sich.}
\end{center}

\textbf{Auslöschung} (catastrophic cancellation): subtrahiert man zwei fast gleiche Zahlen, löschen
sich die führenden Stellen weg, übrig bleibt Rauschen. Gegenmittel: \emph{umformen}, sodass addiert
statt subtrahiert wird, z.\,B.\ $\sqrt{x+1}-\sqrt{x}=\frac{1}{\sqrt{x+1}+\sqrt{x}}$.
Festkomma dagegen hat \emph{konstante absolute} Auflösung $2^{-f}$ (gleichmäßiges Raster), dafür
kleinen Wertebereich.

\begin{rechnung}{Maschinenepsilon \& warum eine Differenz „explodiert``}
System mit $t=2$ Mantissenbits $\Rightarrow \varepsilon=2^{-2}=0{,}25$. Darstellbar ab $1$:
$1{,}00;\,1{,}25;\,1{,}50;\,1{,}75$. Was ist $(1+0{,}1)-1$? Da $0{,}1$ kleiner als die halbe
Auflösung $0{,}125$ ist, rundet $1{,}1\to1{,}0$. Also $(1+0{,}1)-1 = 1{,}0-1{,}0 = 0$ -- der ganze
Beitrag $0{,}1$ ist weg. Eine anschließende Division $1/0$ liefert $\infty$. Lehre: kleine Zahlen
\emph{zuerst} zusammenfassen, kritische Differenzen umformen.
\end{rechnung}

% =====================================================================
\section{Fehleranalyse: die vier Eigenschaften}
\teils

\subsection{Die Begriffe (wer wird mit wem verglichen?)}
Notation: $f$ = wahres Problem, $\hat f$ = Algorithmus, $\tilde x$ = gestörte Eingabe.
\begin{itemize}
\item \textbf{Kondition} $\kappa=|f(x)-f(\tilde x)|$: Empfindlichkeit des \emph{Problems} gegenüber
  Eingabestörung. \emph{Gehört zum Problem.} Faustformel $\kappa\approx|f'(x)|\cdot\delta x$.
\item \textbf{Stabilität} $|\hat f(\tilde x)-\hat f(x)|\le s\,|\tilde x-x|$: verstärkt der
  \emph{Algorithmus} Störungen? \emph{Gehört zum Verfahren.}
\item \textbf{Konsistenz} $c=|\hat f(x)-f(x)|$: wie gut nähert die Diskretisierung bei
  \emph{exakter} Eingabe? Skaliert typischerweise mit $h$.
\item \textbf{Konvergenz} $|\hat f(\tilde x)-f(x)|$: der \emph{Gesamtfehler}.
\end{itemize}
\textbf{Merksatz:} Konsistenz $+$ Stabilität $\Rightarrow$ Konvergenz.

Die Kondition versteht man am besten geometrisch: an einer \emph{flachen} Stelle erzeugt eine
Eingabestörung $\delta x$ nur eine kleine Ausgangsänderung (gut konditioniert), an einer
\emph{steilen} Stelle eine große (schlecht konditioniert):

\begin{center}
\begin{tikzpicture}[scale=0.95]
  \draw[ax](0,0)--(3.1,0) node[right]{$x$};
  \draw[ax](0,0)--(0,3.4) node[above]{$f(x)$};
  \draw[crv,domain=0:2.65,samples=80,smooth] plot(\x,{0.45*\x*\x});
  % flach [0.4,0.9]: f = 0.072 .. 0.3645
  \draw[green,line width=1.2pt](0.4,0)--(0.9,0);
  \draw[green,dashed](0.4,0)--(0.4,{0.45*0.16})--(0,{0.45*0.16});
  \draw[green,dashed](0.9,0)--(0.9,{0.45*0.81})--(0,{0.45*0.81});
  \draw[{Stealth[length=4pt]}-{Stealth[length=4pt]},green](0.12,{0.45*0.16})--(0.12,{0.45*0.81});
  \node[green,right] at (0.15,{0.45*0.5}){\scriptsize $\Delta f$ klein};
  \node[green,below] at (0.65,-0.02){\scriptsize $\delta x$};
  % steil [2.0,2.5]: f = 1.8 .. 2.8125
  \draw[rose,line width=1.2pt](2.0,0)--(2.5,0);
  \draw[rose,dashed](2.0,0)--(2.0,{0.45*4})--(0,{0.45*4});
  \draw[rose,dashed](2.5,0)--(2.5,{0.45*6.25})--(0,{0.45*6.25});
  \draw[{Stealth[length=4pt]}-{Stealth[length=4pt]},rose](0.12,{0.45*4})--(0.12,{0.45*6.25});
  \node[rose,right] at (0.15,{0.45*5.1}){\scriptsize $\Delta f$ groß};
  \node[rose,below] at (2.25,-0.02){\scriptsize $\delta x$};
\end{tikzpicture}\\[-2pt]
{\footnotesize\color{muted}Gleiches $\delta x$: flach $\to$ kleines $\Delta f$ (gut konditioniert), steil $\to$ großes $\Delta f$ (schlecht).}
\end{center}

\begin{rechnung}{Konsistenz auf zwei Gittern}
$g(x)=(x-2)^2$, am nächsten Gitterpunkt ausgewertet. An $x=2{,}7$ (echt $g=0{,}49$):
grobes Gitter $h=1$ rundet auf $3$, $\hat g=g(3)=1\Rightarrow c=0{,}51$;
feines Gitter $h=\tfrac12$ rundet auf $2{,}5$, $\hat g=g(2{,}5)=0{,}25\Rightarrow c=0{,}24$.
Kleineres $h$ ist konsistenter. \textbf{Aber}: unter einem optimalen $h$ dominieren Rundungsfehler
-- kleiner ist \emph{nicht} automatisch besser.
\end{rechnung}

% =====================================================================
\section{Newton- \& Sekantenverfahren}
\teils

\subsection{Was \& warum}
Gesucht ist eine \textbf{Nullstelle} $f(x)=0$, wenn keine Formel existiert. Newton nutzt die
\emph{Tangente}: man ersetzt $f$ lokal durch ihre Linearisierung und springt zu deren Nullstelle.

\subsection{Herleitung aus Taylor (musst du können)}
Taylor 1.\ Ordnung um $x_n$: $f(x)\approx f(x_n)+f'(x_n)(x-x_n)$. Diese Tangente $=0$ setzen und nach
$x$ auflösen:
\[
  0=f(x_n)+f'(x_n)(x-x_n)\ \Longrightarrow\
  \boxed{\,x_{n+1}=x_n-\dfrac{f(x_n)}{f'(x_n)}\,}
\]
Geometrisch: lege die Tangente an, lies ab, wo sie die $x$-Achse schneidet, wiederhole:

\begin{center}
\begin{tikzpicture}[scale=0.9]
  \draw[ax](-0.2,0)--(4.0,0) node[right]{$x$};
  \draw[ax](0,-1.4)--(0,3.3) node[above]{$f(x)$};
  % f = 0.5 x^2 - 2, Wurzel bei x=2; f'(x)=x
  \draw[crv,domain=1.2:3.05,samples=70,smooth] plot(\x,{0.5*\x*\x-2});
  \fill[rose](2,0) circle(1.7pt);
  \node[rose,below right] at (2.0,-0.05){\scriptsize Wurzel ($\sqrt4$)};
  % x0 = 3 : f(3)=2.5, f'(3)=3, Tangente y=2.5+3(x-3) trifft x-Achse bei x1=2.1667
  \fill[yellow](3,2.5) circle(1.7pt);
  \draw[yellow,dashed](3,0) node[below]{\scriptsize $x_0$}--(3,2.5);
  \draw[green,line width=1pt](3,2.5)--(2.1667,0);
  \fill[green](2.1667,0) circle(1.5pt);
  \node[green,above right] at (2.17,0.04){\scriptsize $x_1$};
  \node[green] at (3.55,1.25){\scriptsize Tangente};
\end{tikzpicture}\\[-2pt]
{\footnotesize\color{muted}Tangente in $x_0$ $\to$ ihre Nullstelle $=x_1$ (näher an der Wurzel). Wiederholen.}
\end{center}

\subsection{Konvergenz, Kriterium, Versagen}
\textbf{Quadratische Konvergenz} $|e_{n+1}|\le C|e_n|^2$: die Zahl korrekter Stellen \emph{verdoppelt}
sich pro Schritt. \textbf{Kriterium}: $x_0$ nah genug an einer \emph{einfachen} Wurzel, $f\in C^2$,
$f'\neq0$ dort. \textbf{Versagen}: $f'(x_n)=0$ (Division durch $0$), schlechter Start (Divergenz),
oder ein Zyklus. \textbf{Sekantenverfahren}: ersetzt $f'$ durch die Vorwärtsdifferenz
$\frac{f(x_0+h)-f(x_0)}{h}$ (Fehler $\mathcal O(h)$) bzw.\ durch die Steigung der letzten zwei Punkte
-- ableitungsfrei, dafür langsamer (Ordnung $\approx1{,}618$) und zwei Startwerte nötig.

\begin{rechnung}{Newton für $\sqrt2$ ($f=x^2-2$, $f'=2x$, $x_0=2$)}
\[
x_1=2-\frac{2^2-2}{2\cdot2}=2-\tfrac12=\tfrac32,\qquad
x_2=\tfrac32-\frac{(\tfrac32)^2-2}{2\cdot\tfrac32}=\tfrac32-\frac{1/4}{\mathbf 3}=\tfrac{17}{12}\approx1{,}4167.
\]
Wichtig: $f'(\tfrac32)=2\cdot\tfrac32=\mathbf3$ im Nenner. ($\sqrt2\approx1{,}4142$ -- nach 2 Schritten 3 Stellen.)
\end{rechnung}

% =====================================================================
\section{Optimierung}
\stark

\subsection{Was \& warum}
Gesucht ist ein \textbf{Minimum} von $f$. Ein \emph{lokales} Minimum ist der kleinste Wert in einer
Umgebung, das \emph{globale} der kleinste überhaupt. Lokale Verfahren (Gradientenabstieg, Newton)
folgen nur lokaler Steigung und laufen ins \emph{nächste} Tal -- bei vielen Tälern verfehlen sie das
globale:

\begin{center}
\begin{tikzpicture}[scale=0.92]
  \draw[ax](0,0)--(8.0,0) node[right]{$x$};
  \draw[ax](0,0)--(0,2.9) node[above]{$f(x)$};
  \draw[crv,domain=0.5:7.5,samples=160,smooth]
     plot(\x,{0.5*sin(1.6*\x r)+0.07*(\x-4)*(\x-4)+1.2});
  % echte Minima: global x≈3.05, lokal x≈6.57
  \fill[accentlt](3.05,{0.5*sin(1.6*3.05 r)+0.07*(3.05-4)*(3.05-4)+1.2}) circle(2pt);
  \node[accentlt,below=2pt] at (3.05,{0.5*sin(1.6*3.05 r)+0.07*(3.05-4)*(3.05-4)+1.2}){\scriptsize \textbf{global}};
  \fill[rose](6.57,{0.5*sin(1.6*6.57 r)+0.07*(6.57-4)*(6.57-4)+1.2}) circle(2pt);
  \node[rose,below=2pt] at (6.57,{0.5*sin(1.6*6.57 r)+0.07*(6.57-4)*(6.57-4)+1.2}){\scriptsize lokal};
  % Ball rollt von x=5.5 in das lokale (schlechtere) Tal
  \fill[yellow](5.5,{0.5*sin(1.6*5.5 r)+0.07*(5.5-4)*(5.5-4)+1.2}) circle(2.3pt);
  \node[yellow,above=1pt] at (5.5,{0.5*sin(1.6*5.5 r)+0.07*(5.5-4)*(5.5-4)+1.2}){\scriptsize Start};
  \draw[yellow,-{Stealth[length=4pt]}](5.75,1.5) to[bend left=20] (6.45,1.28);
\end{tikzpicture}\\[-2pt]
{\footnotesize\color{muted}Ball rollt ins \emph{nächste} (lokale) Tal $x\!\approx\!6{,}6$ -- verfehlt das tiefere globale bei $x\!\approx\!3{,}1$.}
\end{center}

\subsection{Verfahren}
\textbf{Gradientenabstieg} $x_{n+1}=x_n-\eta f'(x_n)$ (Schritt bergab; $\eta$ zu groß $\to$ Oszillation,
zu klein $\to$ langsam). \textbf{Newton fürs Optimum} $x_{n+1}=x_n-\frac{f'(x_n)}{|f''(x_n)|}$ -- der
Betrag $|f''|$ erzwingt einen Schritt \emph{immer bergab} (sonst liefe man bei negativer Krümmung zu
einem Maximum); Klassifikation per $f''$: $>0$ Minimum, $<0$ Maximum. \textbf{Global}: Random Restarts
($P=1-(1-p)^K$), Simulated Annealing (akzeptiert Verschlechterung mit $e^{-\Delta E/T}$, Abkühlen
$T\downarrow$), Basin Hopping. Bei hochfrequenten Störtermen hilft \emph{Glätten}: Schrittweite
$h\approx$ Periode des Störterms wählen, dann mittelt sich die Oszillation heraus.

\begin{rechnung}{Konvexität \& Random Restarts}
Ist $f$ \textbf{konvex} ($f''\ge0$ überall), gibt es nur \emph{ein} Minimum $\Rightarrow$ jedes lokale
ist global, lokale Verfahren genügen. Sonst: $K$ zufällige Starts, Trefferchance $p=\tfrac12$ pro
Start. $P(\ge1\text{ Treffer})=1-(1-\tfrac12)^K$. Für $P>0{,}9$: $(\tfrac12)^K<0{,}1\Rightarrow 2^K>10
\Rightarrow K=4$ ($P=\tfrac{15}{16}$).
\end{rechnung}

% =====================================================================
\section{Differenzialgleichungen (Euler \& Runge-Kutta)}
\offen

\subsection{Was \& warum}
Ein Anfangswertproblem $y'=f(t,y),\ y(t_0)=y_0$ hat oft keine geschlossene Lösung. Man läuft in
Schritten der Weite $h$ voran und schätzt die nächste Stelle aus der Steigung $f$.
\textbf{Explizites Euler}: $y_{n+1}=y_n+h\,f(t_n,y_n)$ -- billig, aber nur für kleine $h$ stabil
(für $y'=-\lambda y$: stabil nur bei $0<h<\tfrac2\lambda$). \textbf{Implizit} nutzt $f(t_{n+1},y_{n+1})$
-- teurer (Gleichung lösen), dafür stabil auch bei großem $h$ (wichtig bei \emph{steifen} Problemen).

\begin{center}
\begin{tikzpicture}[scale=0.92]
  \draw[ax](0,0)--(5.0,0) node[right]{$t$};
  \draw[ax](0,0)--(0,3.4) node[above]{$y$};
  % exakte Lösung y' = -0.5 y, y0=3 -> 3 e^{-0.5 t}
  \draw[crv,domain=0:4.4,samples=80,smooth] plot(\x,{3*exp(-0.5*\x)});
  \node[accent] at (3.7,1.1){\scriptsize exakt $3e^{-0{,}5t}$};
  % Euler h=1: 3, 1.5, 0.75, 0.375, 0.1875
  \draw[rose,very thick] (0,3)--(1,1.5)--(2,0.75)--(3,0.375)--(4,0.1875);
  \foreach \p in {(0,3),(1,1.5),(2,0.75),(3,0.375),(4,0.1875)} \fill[rose]\p circle(1.5pt);
  \node[rose] at (2.15,1.35){\scriptsize Euler ($h{=}1$)};
  \foreach \x in {1,2,3,4} \draw[rulecol](\x,-0.07)--(\x,0.07) node[below]{\scriptsize \x};
\end{tikzpicture}\\[-2pt]
{\footnotesize\color{muted}Euler (rose) folgt geraden Tangentenstücken und liegt unter der Kurve -- RK4 wertet $f$ viermal pro Schritt aus und trifft fast exakt.}
\end{center}

\subsection{Runge-Kutta \& Ordnung}
\textbf{RK4} kombiniert vier Steigungen: $y_{n+1}=y_n+\frac h6(k_1+2k_2+2k_3+k_4)$ mit
$k_1=f(t_n,y_n)$, $k_2=f(t_n{+}\tfrac h2,y_n{+}\tfrac h2k_1)$, $k_3=f(t_n{+}\tfrac h2,y_n{+}\tfrac
h2k_2)$, $k_4=f(t_n{+}h,y_n{+}hk_3)$. \textbf{Konvergenzordnung} $p$: globaler Fehler $\mathcal O(h^p)$
-- Euler $p{=}1$, RK4 $p{=}4$ (darum bei gleichem $h$ enorm genauer). Ein \textbf{Butcher-Tableau}
$\begin{smallmatrix}c&|&A\\\hline&|&b^{\top}\end{smallmatrix}$ kodiert Knoten $c_i$, Gewichte $b_i$,
Kopplung $a_{ij}$; \emph{explizit} $\Leftrightarrow$ $A$ streng untere Dreiecksmatrix.

\begin{rechnung}{Ein Euler- vs.\ ein RK4-Schritt ($y'=-y$, $y_0=1$, $h=1$, exakt $e^{-1}=0{,}368$)}
Euler: $y_1=1+1\cdot(-1)=0$. \quad RK4: $k_1{=}-1,\ k_2{=}f(1{-}\tfrac12){=}-0{,}5,\
k_3{=}f(1{-}0{,}25){=}-0{,}75,\ k_4{=}f(1{-}0{,}75){=}-0{,}25$,
$y_1=1+\tfrac16(-1-1-1{,}5-0{,}25)=0{,}375$. RK4 trifft fast exakt, Euler weit daneben.
\end{rechnung}

% =====================================================================
\section{Integration \& Interpolation}
\ueben

\subsection{Numerische Integration (Newton-Cotes)}
Idee: $\int_a^b f\,dx$ approximieren, indem man $f$ durch ein einfaches Polynom über Stützpunkte
ersetzt und \emph{das} integriert. Gerade $\to$ \textbf{Trapez}, Parabel $\to$ \textbf{Simpson}. Der
\emph{Exaktheitsgrad} ist der höchste Polynomgrad, der noch fehlerfrei integriert wird (Trapez $1$,
Simpson $\mathbf 3$). Geometrisch ersetzt Trapez die Kurve durch eine Sehne -- der Spalt dazwischen
ist der Fehler; Simpson legt eine Parabel an und trifft Krümmung viel besser:

\begin{center}
\begin{tikzpicture}[scale=0.95]
  \draw[ax](0,0)--(5.4,0) node[right]{$x$};
  \draw[ax](0,0)--(0,3.2) node[above]{$f$};
  \draw[crv,domain=0.3:5,samples=80,smooth] plot(\x,{0.6+0.12*(\x-2.5)*(\x-2.5)+0.2*\x});
  % Trapez-Sehne von a=1 bis b=4
  \def\fa{0.6+0.12*(1-2.5)^2+0.2*1}   % a
  \def\fb{0.6+0.12*(4-2.5)^2+0.2*4}   % b
  \fill[rose,opacity=0.18] (1,0)--(1,{\fa})--(4,{\fb})--(4,0)--cycle;
  \draw[rose,very thick](1,{\fa})--(4,{\fb});
  \draw[rulecol,dashed](1,0)node[below]{\scriptsize $a$}--(1,{\fa});
  \draw[rulecol,dashed](4,0)node[below]{\scriptsize $b$}--(4,{\fb});
  \node[rose] at (2.5,0.55){\scriptsize Sehne (Trapez)};
  \node[accent] at (3.9,2.7){\scriptsize echte Kurve};
\end{tikzpicture}\\[-2pt]
{\footnotesize\color{muted}Der Spalt zwischen Sehne und Kurve ist der Trapez-Fehler -- bei Krümmung groß.}
\end{center}

\textbf{Zusammengesetzte Gewichte} (für $[0,2]$ mit $f_0,f_1,f_2$, $h=1$):
\[
  \text{Trapez: } \tfrac h2\bigl[f_0+\mathbf2 f_1+f_2\bigr],\qquad
  \text{Simpson: } \tfrac h3\bigl[f_0+\mathbf4 f_1+f_2\bigr]\ \text{\small(„1,4,1 durch 3``)}.
\]
\textbf{Genau ein} innerer Term, mit Gewicht $2$ bzw.\ $4$ -- nicht $2(f_1{+}f_2)$, keine Extra-Terme.
Weiter: \textbf{$3/8$-Regel} (Gewichte $1,3,3,1$), \textbf{Boole} ($7,32,12,32,7$),
\textbf{Romberg} (Trapez für $h,\tfrac h2$ extrapolieren: $R=\frac{4T(h/2)-T(h)}{3}$),
\textbf{Monte-Carlo} ($\hat I=|\Omega|\frac1N\sum f(X_i)$, Fehler $\propto 1/\sqrt N$, dimensions\-unabhängig
$\Rightarrow$ schlägt das Gitter ($n^d$ Punkte) in hohen Dimensionen).

\subsection{Interpolation}
Durch $n{+}1$ Punkte gibt es \emph{genau ein} Polynom vom Grad $\le n$. \textbf{Lagrange}:
$p(x)=\sum_i y_i L_i(x)$ mit $L_i(x)=\prod_{j\neq i}\frac{x-x_j}{x_i-x_j}$ ($L_i$ ist $1$ am eigenen
Knoten, sonst $0$). \textbf{Vorsicht hoher Grad}: auf gleichmäßigen Knoten oszilliert das Polynom an
den Rändern stark -- das \textbf{Runge-Phänomen}. Darum nicht den Grad hochdrehen, sondern
\emph{zusammengesetzte} Regeln oder \textbf{Splines} (stückweise niedriger Grad, glatt):

\begin{center}
\begin{tikzpicture}[scale=0.9]
  \draw[ax](-3.0,0)--(3.2,0) node[right]{$x$};
  \draw[ax](0,0)--(0,2.7) node[above]{$y$};
  \draw[green,very thick,domain=-2.7:2.7,samples=100,smooth] plot(\x,{2/(1+\x*\x)});
  \node[green] at (1.75,1.5){\scriptsize wahre Funktion};
  % hochgradiges Polynom: trifft Knoten ±1,±2, schwingt an Rändern über
  \draw[rose,very thick,domain=-2.7:2.7,samples=160,smooth]
     plot(\x,{2/(1+\x*\x) + 0.06*(\x*\x-1)*(\x*\x-4)});
  \node[rose] at (-1.9,2.1){\scriptsize hoher Grad};
  \foreach \x in {-2,-1,1,2} \fill[accentlt](\x,{2/(1+\x*\x)})circle(1.7pt);
\end{tikzpicture}\\[-2pt]
{\footnotesize\color{muted}Runge: das Polynom (rose) trifft die Knoten ($\bullet$), schwingt aber an den Rändern stark über -- Splines vermeiden das.}
\end{center}

\begin{rechnung}{Lagrange durch $(0,1),(1,2),(2,5)$}
$L_0=\frac{(x-1)(x-2)}{(0-1)(0-2)}=\frac{(x-1)(x-2)}2$,\;
$L_1=\frac{x(x-2)}{(1)(-1)}=-x(x-2)$,\;
$L_2=\frac{x(x-1)}{(2)(1)}=\frac{x(x-1)}2$. Dann
$p=1\!\cdot\!L_0+2\!\cdot\!L_1+5\!\cdot\!L_2=x^2+1$ (Probe: $1,2,5$ \checkmark), $p(\tfrac12)=\tfrac54$.
\end{rechnung}

% =====================================================================
\section{Fourier-Analyse \& FFT}
\offen

\subsection{Was \& warum}
Die Fourier-Transformation $F(\omega)=\int f(t)e^{-i\omega t}\,dt$ zerlegt ein Zeitsignal in seine
\textbf{Frequenzen}: $|F(\omega)|$ sagt, wie stark Frequenz $\omega$ enthalten ist; die
Rücktransformation baut das Signal wieder zusammen. Mit der Euler-Formel
$\cos(\omega t)=\tfrac12(e^{i\omega t}+e^{-i\omega t})$ liefert jeder reelle Kosinus \emph{zwei}
Linien bei $\pm\omega$. Ein Summensignal („Interferenz``) zeigt also pro Schwingung einen Peak:

\begin{center}
\begin{tikzpicture}[scale=0.9]
  % links: Zeitsignal sin(2t) + 1/2 sin(4t)  (Amplituden 1 und 1/2)
  \draw[ax](0,-1.7)--(0,1.7);
  \draw[ax](-0.1,0)--(4.4,0) node[right]{\scriptsize $t$};
  \draw[crv,domain=0:4,samples=160,smooth] plot(\x,{sin(2*\x r)+0.5*sin(4*\x r)});
  % Pfeil
  \draw[muted,-{Stealth[length=5pt]}](4.8,0)--(5.8,0) node[midway,above]{\scriptsize FFT};
  % rechts: Spektrum -- Höhen 2:1 wie die Amplituden, Peaks bei w1 und 2 w1
  \begin{scope}[xshift=6.3cm]
   \draw[ax](0,0)--(3.6,0) node[right]{\scriptsize $\omega$};
   \draw[ax](0,0)--(0,1.75) node[above]{\scriptsize $|F|$};
   \draw[rose,line width=1.4pt](1.2,0)--(1.2,1.3); \fill[rose](1.2,1.3)circle(1.6pt);
   \draw[rose,line width=1.4pt](2.4,0)--(2.4,0.65); \fill[rose](2.4,0.65)circle(1.6pt);
   \node[rose,below=1pt] at (1.2,0){\scriptsize $\omega_1$};
   \node[rose,below=1pt] at (2.4,0){\scriptsize $2\omega_1$};
  \end{scope}
\end{tikzpicture}\\[-1pt]
{\footnotesize\color{muted}Summensignal (links) $\to$ Spektrum (rechts): zwei Peaks bei $\omega_1,2\omega_1$, Höhen $1:\tfrac12$ wie die Amplituden.}
\end{center}

\textbf{DFT}: $X_k=\sum_{n=0}^{N-1}x_n e^{-i2\pi kn/N}$ ($N$ diskrete Bins), Kosten $\mathcal O(N^2)$.
Die \textbf{FFT} spaltet rekursiv in gerade/ungerade Indizes $\Rightarrow$ $\mathcal O(N\log N)$.
Bin $k$ entspricht Frequenz $f_k=k\,\frac{f_s}{N}$; die \textbf{Nyquist}-Frequenz $f_s/2$ ist die
höchste eindeutig darstellbare -- höhere erscheinen als \emph{Aliasing} verfälscht.

\begin{rechnung}{DFT von $x=(1,0,1,0)$, $N=4$}
$X_0=1{+}0{+}1{+}0=2$;\; $X_1=1+1\cdot e^{-i\pi}=1-1=0$;\;
$X_2=1+1\cdot e^{-i2\pi}=2$;\; $X_3=0$. Also $|X|=(2,0,2,0)$: Gleichanteil ($X_0$) plus
schnellste Schwingung ($X_2$, Nyquist) -- genau die Struktur von $(1,0,1,0)$.
\end{rechnung}

% =====================================================================
\section{Modelltraining (Numerik trifft ML)}
\ueben

\subsection{Was \& warum}
Ein lineares Modell $\hat y(x)=wx+b$ wird trainiert, indem ein \textbf{Verlust} minimiert wird --
also Optimierung (Abschnitt 4) auf einer Verlustfunktion. Standard ist der \textbf{mittlere
quadratische Fehler}
\[
  L(w,b)=\frac1n\sum_{i=1}^n\bigl(wx_i+b-y_i\bigr)^2 .
\]
Die Abweichungen werden \textbf{quadriert} -- sonst heben sich $+$/$-$ auf; das Quadrat bestraft große
Fehler stärker und macht $L$ glatt/konvex. Trainiert wird per Gradientenabstieg.

\subsection{Zwei Dinge, die immer drankommen}
\textbf{(1) Erster Loss = Varianz.} Mit $w=0,\ b=\bar y$ sagt das Modell überall den Mittelwert
voraus, also $L=\frac1n\sum(\bar y-y_i)^2=\mathrm{Var}(y)$ -- die Varianz ist das Mittel der
\emph{quadrierten} Abweichungen. \textbf{(2) Numerik-Brille auf ML-Praktiken:}
\begin{itemize}
\item Überparametrisiertes Modell soll eine kleine Teilmenge auf Loss $0$ bringen -- klappt das nicht,
  ist es ein \emph{Bug} (Modell/Loss/Daten), nicht der Optimierer (wie „Solver an bekannter Lösung testen``).
\item Mini-Batch-SGD $\leftrightarrow$ Monte-Carlo (Gradient als Stichprobenmittel, Fehler $\propto1/\sqrt B$).
\item Mehrere Seeds, bester Lauf $\leftrightarrow$ Random Restarts.
\item Quantisierung $\leftrightarrow$ Präzisionswahl im Gleitkomma (kleinstes ausreichendes Format).
\end{itemize}

\begin{rechnung}{Varianz nicht ohne Quadrat!}
$y=(2,4,6)$, $\bar y=4$. Falsch: $\tfrac13(-2+0+2)=0$. Richtig:
$\mathrm{Var}=\tfrac13[(-2)^2+0^2+2^2]=\tfrac13(4+0+4)=\tfrac83\approx2{,}67$. Das ist der erwartete
erste Loss; ein gedruckter Wert $\approx100$ verrät einen Bug/falsche Initialisierung.
\end{rechnung}

\vspace{4pt}
\noindent\textcolor{rulecol}{\rule{\linewidth}{0.4pt}}\\[2pt]
{\small\color{muted}Stand-Legende: \stark\ Themen 1,4 $\cdot$ \teils\ Themen 2,3 $\cdot$
\ueben\ Themen 6,8 $\cdot$ \offen\ Themen 5,7. Letztere zwei kamen in der letzten Klausur nicht dran,
gehören aber zum Stoff -- für die nächste Klausur mitnehmen.}

\end{document}
